% =============================================================================
\section{The Panoramic Data}
\label{sec:panorama}
% =============================================================================

Every track in \cref{fig:pipeline} consumes the same three artifacts: a
$360^\circ$ panorama, an aligned metric depth map, and a per-pixel
semantic typing. The depth map and semantic typing are both derived from
the generated panorama, which extends the provided input image into a complete
$360^\circ$ view of the scene.

We use one capture, a photograph of Mount Rainier, as a running example throughout.
\Cref{fig:rainier-input} shows the input and the panorama generated from it: the
photograph covers roughly a sixth of the sphere, and everything else is synthesized.

\begin{figure}[htbp]
  \centering
  \begin{subfigure}{0.48\linewidth}
    \centering
    \includegraphics[width=\linewidth]{rainier/input.jpg}
    \caption{Input photograph.}
    \label{fig:rainier-photo}
  \end{subfigure}\hfill
  \begin{subfigure}{0.48\linewidth}
    \centering
    \includegraphics[width=\linewidth]{rainier/depth-perspective.png}
    \caption{Perspective depth, the metric anchor.}
    \label{fig:rainier-depth-persp}
  \end{subfigure}

  \vspace{0.6em}
  \includegraphics[width=\linewidth]{rainier/panorama.jpg}
  \caption{The Mount Rainier example. The input photograph (a) and its perspective
  depth estimate (b) are outpainted into a full equirectangular panorama (below).
  The photograph occupies the region around the mountain; the remaining
  $\sim$\SI{83}{\percent} of the sphere is generated.}
  \label{fig:rainier-input}
\end{figure}

\subsection{Panorama Generation and Metric Calibration}

The input image is first processed by generating a synthetic caption and a
perspective depth estimate, then outpainting the image to a full equirectangular
panorama, which is super-sampled to its working resolution. Panoramic depth is
estimated on the result~\cite{lin2025dap}.

\emph{Depth Any Panoramas} is built for equirectangular imagery, but its output is
relative rather than metric. We calibrate it against
Depth Anything~3~\cite{lin2025depth3recoveringvisual} run on the original input,
which gives a metric baseline over the region the photograph actually observed.
The two are related by a monotonic response curve rather than a single scale
factor, because the relationship is not linear.

\paragraph{Fitting the response curve.}
Let $r \in [0,1]$ be the panoramic model's raw prediction and $d$ the metric depth
from the perspective model. We project the hero photograph into the panorama using
the same field of view that was used to place it during outpainting, which gives a
set of corresponding samples $\{(r_i, d_i)\}$ over the overlap region. Sky pixels
and non-finite or non-positive samples are discarded. The surviving samples are
partitioned into $N = 14$ quantile bins of $r$, and each bin contributes one knot

\begin{equation}
  c_b = \operatorname{median}\{r_i : i \in \mathcal{B}_b\},
  \qquad
  m_b = \operatorname{median}\{d_i : i \in \mathcal{B}_b\},
\end{equation}

where $\mathcal{B}_b$ is the set of samples in bin $b$. Medians rather than means
keep mismatched samples, typically foreground objects surviving into the overlap,
from dragging a knot, and a running maximum $\tilde{m}_b = \max_{k \le b} m_k$
enforces that depth cannot decrease as the raw prediction increases. Calibration is
attempted only above $300$ valid samples; below that the overlap is too small to fit
reliably and the depth is left uncalibrated. Metric depth is then recovered by
interpolating $(c_b, \tilde{m}_b)$. The curve is treated as a property of the model
rather than of the image, so the same fit applies to every panoramic depth map in
the scene, which is what keeps them on a common scale.

\paragraph{Extrapolating beyond the fitted range.}
The overlap region only ever covers the photograph's own near field, a few to a few
tens of meters, so it never contains genuinely distant terrain. Interpolation alone
would flat-clamp every raw value past the fitted range onto the curve's last knot,
collapsing a mountain a kilometer away and a tree just past the fitted range onto
the same depth. We instead extrapolate in log-depth space, since the raw prediction
approaches its ceiling asymptotically as true depth grows,

\begin{equation}
  d(r) = \exp\!\big(\log \tilde{m}_{N} + s\,(r - c_{N})\big),
  \qquad
  s = \frac{\log \tilde{m}_{N} - \log \tilde{m}_{N-k}}{c_{N} - c_{N-k}},
\end{equation}

with the slope $s$ estimated from the last $k = 3$ knots, and symmetrically at the
near end, capped at $d \le \lambda\,\tilde{m}_N$ with $\lambda = 20$. The cap
matters because the raw output saturates well before true distance does: a
mountainside a few hundred meters out and a hard-clamped sky pixel read at
essentially the same ceiling. A slope estimated from a few tail bins and then
exponentiated over that gap is extremely noise-sensitive, and produced fabricated
depths of tens of kilometers for every pixel sharing a saturated value. One
outlier of that size dominates anything treating panorama depth as a trustworthy
scale, silently discarding real terrain as out of bounds rather than merely far.
The cap recovers no distance the model cannot resolve, but it keeps far content
bounded and ordered.

\paragraph{Fitting the scene into the terrain budget.}
The calibrated map must finally fit inside the terrain grid's footprint,
$D_{\max} = \SI{100}{\meter}$, half the \SI{200}{\meter} grid extent. A uniform
rescale preserves relative distances exactly and is the wrong invariant to protect:
the scale would be set by the most distant thing in the scene, effectively unbounded
for a landscape, and it invalidates the pipeline's second and more load-bearing
anchor, the assumed camera height that places ground at $-h$ below the camera. On
Rainier, fitting a summit kilometers away into \SI{100}{\meter} divided the map by
$6.5\times$: ground \SI{1.37}{\meter} below the camera returned \SI{0.21}{\meter} of
depth, putting the reconstructed surface \SI{0.9}{\meter} \emph{above} eye level
with the camera buried in its own terrain. A flat clamp fails oppositely, collapsing
a \SI{120}{\meter} ridge and a \SI{900}{\meter} peak onto one wall.

We compress only the far end, leaving the near field metrically intact. With a knee
at $k = 0.6\,D_{\max}$ and tail $t = D_{\max} - k$,

\begin{equation}
  d' =
  \begin{cases}
    d, & d \le k, \\[2pt]
    k + t\left(1 - \exp\!\left(-\dfrac{d - k}{t}\right)\right), & d > k.
  \end{cases}
\end{equation}

This is $C^1$ at the knee, strictly monotonic, and asymptotic to $D_{\max}$, so no
pixel exceeds the budget without a hard clamp. Distant content keeps its ordering;
it simply stops being proportional, which it never meaningfully was once a mountain
kilometers away had to fit a \SI{200}{\meter} grid. Because the curve is fixed by
$(D_{\max}, k)$ rather than by the data's maximum, saturated sky cannot contaminate
it, and applying the same function to both depth maps places the same point at the
same distance in each.

\begin{figure}[htbp]
  \centering
  \includegraphics[width=\linewidth]{rainier/depth-calibrated.png}
  \caption{Calibrated panoramic depth for the Rainier capture, near field bright.
  Two properties discussed above are visible directly. The near-field meadow is
  resolved in detail, while everything from the mid-ground treeline outward
  collapses into a narrow band of values: the raw prediction has saturated, and the
  far-range compression maps that whole range onto the approach to
  $D_{\max}$. The sky is masked out rather than assigned a depth.}
  \label{fig:rainier-depth}
\end{figure}

\paragraph{Consequences.}
The result is metrically compressed: near distances are accurate, far ones
systematically shrunk. Since the user explores only a few meters of the scene the
visual effect is minimal, relying on forced perspective, but it is not free.
\Cref{sec:results} returns to a defect it causes, where object size is computed as
angular extent times depth and inherits the same shrinkage.

\subsection{Object-Removed Panorama}

To assist terrain generation we produce a second panorama with foreground objects
removed, using ObjectClear~\cite{zhao2026objectclear}, which removes objects
together with their effects such as reflections and cast shadows. Objects are
selected by depth threshold, taking those closest to the camera.

Depth taken from the panorama with objects still in it produces spurious values and
raised bumps wherever an object stands in front of the ground, since the surface
behind it was never observed; removing them first gives the depth model an
unobstructed view. \Cref{fig:rainier-object-removed} shows the effect: the near-field
meadow dominating the lower half is replaced by the bare ground the terrain track
needs. The edit is destructive by design. The flowers are not discarded but
reconstructed separately as a ground-cover population
(\cref{sec:distribution}), and the region typing that places them is read from the
\emph{original} panorama.

\begin{figure}[htbp]
  \centering
  \includegraphics[width=\linewidth]{rainier/panorama.jpg}\\[0.4em]
  \includegraphics[width=\linewidth]{rainier/panorama-object-removed.jpg}
  \caption{Foreground removal. The original panorama (top) and the object-removed
  panorama (bottom) used for terrain reconstruction. The near-field meadow and the individual
  trees have been removed together with their shadows, exposing the ground surface
  behind them. Depth taken from the top image would place the flower canopy, not the
  ground, at the bottom of the height field.}
  \label{fig:rainier-object-removed}
\end{figure}

\subsection{Semantic Region Map}

SegFormer~\cite{DBLP:journals/corr/abs-2105-15203} assigns per-pixel region types
(sky, water, terrain, ground, vegetation, built objects, road, trail) over the
panorama. This map drives much of the downstream scene understanding, from terrain
construction to object placement.

\begin{figure}[htbp]
  \centering
  \includegraphics[width=\linewidth]{rainier/regions.jpg}
  \caption{Resolved semantic region typing for the Rainier capture, overlaid on the
  panorama. Vegetation, bare terrain, snow and standing water separate cleanly in
  the near field. The typing is run twice, once on the original panorama and once on
  the object-removed one, because the two answer different questions: where ground cover
  should be scattered is a property of the scene as photographed, while what the
  terrain surface is made of is a property of the scene with the objects taken away.}
  \label{fig:rainier-regions}
\end{figure}
